% Part: first-order-logic
% Chapter: beyond
% Section: modal-logics

\documentclass[../../../include/open-logic-section]{subfiles}

\begin{document}

\olfileid[pt]{fol}{byd}{mod}

\olsection{Lógicas Modais}

Considere o seguinte exemplo de uma sentença condicional:
\begin{quote}
  Se Jeremy está sozinho naquela sala, então ele está bêbado, nu e
  dançando nas cadeiras.
\end{quote}
Este é um exemplo de uma asserção condicional que pode ser
materialmente verdadeira, mas ainda assim enganosa, pois parece sugerir
que há um vínculo mais forte entre o antecedente e a conclusão do que
simplesmente o fato de que ou o antecedente é falso ou a conclusão é
verdadeira. Ou seja, a formulação sugere que a afirmação não é apenas
verdadeira neste mundo particular (onde pode ser trivialmente verdadeira,
porque Jeremy não está sozinho na sala), mas que, além disso, a conclusão
\emph{teria sido} verdadeira \emph{se} o antecedente tivesse sido
verdadeiro. Em outras palavras, pode-se tomar a asserção como significando
que a afirmação é verdadeira não apenas neste mundo, mas em qualquer mundo
``possível''; ou que ela é \emph{necessariamente} verdadeira, em oposição
a ser apenas verdadeira neste mundo particular.

A lógica modal foi concebida para dar sentido a esse tipo de
necessidade. Obtém-se a lógica proposicional modal a partir da
lógica proposicional ordinária adicionando um operador caixa; isto é,
se $!A$ é !!a{formula}, então $\Box !A$ também o é.  Intuitivamente,
$\Box !A$ afirma que $!A$ é \emph{necessariamente} verdadeira, ou
verdadeira em qualquer mundo possível. $\Diamond !A$ é usualmente
tomada como abreviação de $\lnot \Box \lnot !A$, e pode ser lida como
afirmando que $!A$ é \emph{possivelmente} verdadeira. É claro que a
modalidade também pode ser adicionada à lógica de predicados.

!!^{structure}s de Kripke podem ser usadas para fornecer uma semântica
para a lógica modal; de fato, Kripke projetou essa semântica tendo a
lógica modal em mente. Em vez de restringir a ordens parciais, de modo
mais geral tem-se um conjunto de ``mundos possíveis'', $P$, e uma
relação binária de ``acessibilidade'' $\Atom{R}{x,y}$ entre mundos.
Intuitivamente, $\Atom{R}{p,q}$ afirma que o mundo~$q$ é compatível
com~$p$; isto é, se estamos ``em'' o mundo~$p$, temos que considerar a
possibilidade de que o mundo pudesse ter sido como~$q$.

A lógica modal é às vezes chamada de lógica ``intensional'', em oposição
a uma ``extensional''. A semântica pretendida para uma lógica
extensional, como a lógica clássica, refere-se apenas a um único mundo,
o ``atual''; enquanto a semântica para uma lógica ``intensional'' depende
de uma ontologia mais elaborada. Além de estruturar a necessidade, pode-se
usar a modalidade para estruturar outras construções linguísticas,
reinterpretando $\Box$ e $\Diamond$ de acordo com a aplicação. Por exemplo:
\begin{enumerate}
\item Na lógica de demonstrabilidade, $\Box !A$ é lida como ``$!A$ é
  demonstrável'' e $\Diamond !A$ é lida como ``$!A$ é consistente''.
\item Na lógica epistêmica, pode-se ler $\Box !A$ como ``eu sei
  $!A$'' ou ``eu acredito em $!A$''.
\item Na lógica temporal, pode-se ler $\Box !A$ como ``$!A$ é sempre
  verdadeira'' e $\Diamond !A$ como ``$!A$ é às vezes verdadeira''.
\end{enumerate}

Gostaríamos de aumentar a lógica com regras e axiomas que lidem com
modalidade. Por exemplo, o sistema $\Log{S4}$ consiste nos axiomas e
regras ordinários da lógica proposicional, juntamente com os seguintes
axiomas:
\begin{align*}
& \Box (!A \lif !B) \lif (\Box !A \lif \Box !B)\\
& \Box !A \lif !A \\
& \Box !A \lif \Box \Box !A
\intertext{assim como uma regra, ``de $!A$ conclua $\Box !A$.''
  $\Log{S5}$ adiciona o seguinte axioma:}
& \Diamond !A \lif \Box \Diamond !A
\end{align*}
Variações desses axiomas podem ser adequadas para diferentes aplicações;
por exemplo, S5 é usualmente tomado para caracterizar a noção de
necessidade lógica. E o aspecto agradável é que se pode usualmente
encontrar uma semântica para a qual o sistema de !!{derivation} é correto
e completo, restringindo a relação de acessibilidade nas !!{structure}s
de Kripke de maneiras naturais. Por exemplo, $\Log{S4}$ corresponde à
classe de !!{structure}s de Kripke em que a relação de acessibilidade é
reflexiva e transitiva. $\Log{S5}$ corresponde à classe de !!{structure}s
de Kripke em que a relação de acessibilidade é {\em universal}, isto é,
todo mundo é acessível a partir de qualquer outro; assim $\Box !A$ vale
se, e somente se, $!A$ vale em todo mundo.

\end{document}
